\scanpage{98}

\[
\sum|x_i x_i'|\leq N(x)N'(x')\unsure{)}\tag*{(10)}
\]

En effet, on a $\sum_{1}^{n}|x_i x_i'|=\langle |x|_n,|x'|_n\rangle\leq N(|x|_n)N(|x'|_n)\leq N(x)N'(x')$\unsure{)}, d'où (10). Il en résulte un \uline{accouple-}\newline\uline{ment} entre $\ell^N$ et $\ell^{N'}$, \deleted{\unsure{qui} \inserted{\unsure{est}}} \inserted{d'où en restreignant} \uline{entre $\ell_N$ et $\ell^{N'}$}; \deleted{\unsure{cet accoupl} \unsure{on a ainsi associé à $\ell$ de $\ell^{N'}$} \unsure{i.e. une application|une forme} linéaire continue \inserted{\unsure{de norme} $\leq 1$}, $\ell^{N'}$ dans $(\ell_N)'$} \unsure{je dis} que \uline{le dual de $\ell_N$ s'identifie exactement à $\ell^{N'}$} (avec sa norme) par cet accouplement.

\margin{\textcolor{red}{!}\quad\tikz[baseline=-.6ex]{\draw[red,line width=.7pt] (0,0)--(0,.6);}}

En effet, \deleted{\unsure{soit la}} il suffit de montrer qu'une forme linéaire continue $X'$ sur $\ell_N$ est définie par un \inserted{$x'$} élément de $\ell^{N'}$ \deleted{de \illegible{1}} tel que $N'(x')\leq\lVert X'\rVert$.

En effet, $X'$ \deleted{\unsure{est entièrement} \inserted{\unsure{dé|aussi}}} il suffit de le voir pour la \deleted{\unsure{part}} restriction à $\ell^1$, alors $X'$ \inserted{est} défini \deleted{\unsure{p|par}} \inserted{par} un $x'\in\ell^\infty$, \deleted{\unsure{avec}} on a $N'(|x'|_n)\leq\lVert X'\rVert$, d'où aussitôt $N'(x')\leq\lVert X'\rVert$, et le résultat.

-- Notons qu'en général, le dual de $\ell^N$ n'est pas $\ell^{N'}$, \deleted{\unsure{car}} car $\ell_N$ ne sera pas dense dans $\ell^N$. Ex\,: $N=N_\infty$. Notons que

\margin{\uline{Prop}}
\uline{pour que $\ell_N$ soit réflexif, il faut et il suffit que $\ell^N$ resp$^t$ $\ell_{N'}$ resp$^t$ $\ell^{N'}$ le soit}. \inserted{\uline{Il faut et il suffit pour cela que l'on ait $\ell_N=\ell^N$} \uuline{et} \uline{$\ell_{N'}=\ell^{N'}$.}} \deleted{En effet,}

$\ell^N$ réflexif implique $\ell_N$ réflexif (sous-espace fermé) qui implique $\ell^{N'}$ réflexif (dual de $\ell_N$) qui \inserted{de même} implique $\ell_{N'}$ réflexif et lui-$=$ implique $\ell^N$ réflexif, ce qui achève la démonstration \deleted{de l'équivalence pour} circulaire de l'équivalence des 4 conditions. Alors $\ell_{N'}$, qui est un sous-espace \deleted{\unsure{fort}} \uline{total} \inserted{et fortement fermé} du dual $\ell^{N'}$ de $\ell_N$, est identique à $\ell^{N'}$, et pour la $=$ raison on aura $\ell_N=\ell^N$.

\margin{\textcolor{red}{?}}

\inserted{Réciproquement,} s'il en est ainsi, alors le dual de $\ell_N$ est $\ell^{N'}=\ell_{N'}$, et le dual de ce dernier est $\ell^N=\ell_N$, donc $\ell_N$ coïncide avec son bidual, et est donc réflexif.

\inserted{Soit $N$ une norme de Schatten sur $R^{(N)}$ \deleted{non équivalente à $N_\infty$}\inserted{non équivalente à $N_\infty$}.} Soient $E,F$ deux espaces de Hilbert, \deleted{\unsure{$N$|si}} on désigne par $L^{\deleted{\unsure{N}}N}_{\deleted{\unsure{0}}}(E,F)$ le \deleted{\unsure{sous|l'espace}} sous-espace de $L_0(E,F)$ formé des $u$ tels que $\unsure{(\rho_i(u))|(\rho_i u)}\in\ell^N$, i.e. tels que

\[
N(u)=N(\unsure{(\rho_i(u))|(\rho_i u)})<+\infty\tag*{(11)}
\]

\margin{99}
